\section*{Chapter \ref{ch:b1:acetals-protection} --- Carbonyls: Nucleophilic Additions, Acetals and Protection}

\begin{solution}{exo:b1:acetals-protection:1}
\ce{CH3CH(OH)OCH3}: \omterm{def:b1:acetals-protection:hemiacetal}{hemiacetal}. \ce{CH3CH(OCH3)2}: \omterm{def:b1:acetals-protection:hemiacetal}{acetal}.
\ce{CH3CH(OH)2}: hydrate. \ce{CH3OCH3}: ether. \ce{(CH3)2C(OCH2CH3)2}:
\omterm{def:b1:acetals-protection:hemiacetal}{acetal} (of a ketone).
\end{solution}

\begin{solution}{exo:b1:acetals-protection:2}
\ce{CH3CH2CHO + 2CH3OH <=> CH3CH2CH(OCH3)2 + H2O}; propanone with
ethane-1,2-diol gives 2,2-dimethyl-1,3-dioxolane and water,
\ce{(CH3)2CO + HOCH2CH2OH <=> C5H10O2 + H2O}.
\end{solution}

\begin{solution}{exo:b1:acetals-protection:3}
Sodium hydroxide, methylmagnesium bromide and sodium borohydride leave it
unchanged; dilute aqueous sulfuric acid hydrolyses it.
\end{solution}

\begin{solution}{exo:b1:acetals-protection:4}
The OH on carbon 4 adds to the aldehyde carbon 1: a five-membered ring of
four carbons and one oxygen, with an OH on the carbon next to the ring
oxygen.
\end{solution}

\begin{solution}{exo:b1:acetals-protection:5}
Protonation of C=O; methanol adds to carbon; loss of \ce{H+}: \omterm{def:b1:acetals-protection:hemiacetal}{hemiacetal};
protonation of its OH; loss of water to \ce{CH3CH=O+CH3}; second methanol
adds; loss of \ce{H+}: \ce{CH3CH(OCH3)2}. The two protons taken are given
back: the acid is a \omterm{def:b1:elementary-steps:catalyst}{catalyst}.
\end{solution}

\begin{solution}{exo:b1:acetals-protection:6}
Protonation of one OCH3; loss of methanol to \ce{(CH3)2C=O+CH3}; water adds;
loss of \ce{H+}: \omterm{def:b1:acetals-protection:hemiacetal}{hemiacetal}; protonation, loss of methanol, loss of \ce{H+}:
propanone. Excess water keeps $Q$ below $K$ for the hydrolysis, which goes
to completion.
\end{solution}

\begin{solution}{exo:b1:acetals-protection:7}
One water per \omterm{def:b1:acetals-protection:hemiacetal}{acetal}: $0.250 \times 18.0 = \qty{4.5}{g}$, about
\qty{4.5}{mL}. After an hour, $3.2/18.0 = \qty{0.178}{mol}$: \qty{71}{\percent}
conversion.
\end{solution}

\begin{solution}{exo:b1:acetals-protection:8}
With the diol, one molecule replaces two: the reaction does not decrease
the number of molecules (one aldehyde and one diol give one \omterm{def:b1:acetals-protection:hemiacetal}{acetal} and one
water), and the second OH is held close to the reacting carbon, which
closes the ring easily. Both make the cyclic \omterm{def:b1:acetals-protection:hemiacetal}{acetal} more favourable.
\end{solution}

\begin{solution}{exo:b1:acetals-protection:9}
The acid would destroy the \omterm{def:b1:organomagnesium:organometallic}{Grignard reagent} (proton transfer); traces of
water and diol must be removed too.
\end{solution}

\begin{solution}{exo:b1:acetals-protection:10}
1. Protect the ketone of 4-chlorobutan-2-one as a cyclic \omterm{def:b1:acetals-protection:hemiacetal}{acetal}
(ethane-1,2-diol, acid \omterm{def:b1:elementary-steps:catalyst}{catalyst}, Dean--Stark). 2. Magnesium in dry ether:
the \omterm{def:b1:organomagnesium:organometallic}{Grignard reagent} of the protected chloride, no longer destroyed by its
own carbonyl. 3. Add propanone; work-up with aqueous ammonium chloride:
the tertiary alcohol. 4. Dilute aqueous acid: \omterm{def:b1:acetals-protection:protecting-group}{deprotection} to
5-hydroxy-5-methylhexan-2-one.
\end{solution}

\begin{solution}{exo:b1:acetals-protection:11}
$Q = [\text{acetal}][\ce{H2O}]/([\text{aldehyde}][\ce{CH3OH}]^2)$: a large
excess of methanol, squared in the denominator, keeps $Q < K$ until little
aldehyde remains. Removing water acts on the numerator; using methanol as
the solvent is simpler but needs an easily separated product.
\end{solution}

\begin{solution}{exo:b1:acetals-protection:12}
Only the \omterm{def:b1:acetals-protection:hemiacetal}{hemiacetal} OH can leave as water to give a cation stabilised by
the neighbouring ring oxygen (\ce{C=O+}); the other OH groups would have to
leave from ordinary carbons, giving unstabilised cations. Methanol adds
to the stabilised cation: a methyl \omterm{def:b1:acetals-protection:hemiacetal}{acetal}.
\end{solution}

\begin{solution}{pb:b1:acetals-protection:1}
\textbf{1.} \ce{BrMgCH2CH2CH2CHO}.
\textbf{2.} Its carbanion-like carbon would attack the aldehyde of another
molecule (or of itself): \ce{RMgBr} adds to $\mathrm{R'CHO}$, giving an \omterm{def:b1:alcohol-activation:alkoxide}{alkoxide}.
\textbf{3.} The aldehyde, against the \omterm{def:b1:organomagnesium:organometallic}{Grignard reagent}.
\textbf{4.} \omterm{def:b1:acetals-protection:hemiacetal}{Acetals} are stable towards \omterm{def:b1:organomagnesium:organometallic}{Grignard reagents} and bases and are
removed by aqueous acid.
\textbf{5.} It forms more completely (one diol for two alcohol molecules)
and is easy to hydrolyse.
\textbf{6.} $15.09/150.9 = \qty{0.1000}{mol}$.
\textbf{7.} $1.20 \times 0.1000 \times 62.0 = \qty{7.44}{g}$.
\textbf{8.} \ce{BrCH2CH2CH2CHO + HOCH2CH2OH <=> C6H11BrO2 + H2O}
(2-(3-bromopropyl)-1,3-dioxolane).
\textbf{9.} It activates the carbonyl and allows the loss of water in the
second stage; it is regenerated.
\textbf{10.} Toluene and water distil together, condense and fall into the
collector; water sinks and stays, toluene overflows back.
\textbf{11.} $Q$ contains $[\ce{H2O}]$ in its numerator; removing water keeps
$Q < K$, and the reaction proceeds.
\textbf{12.} Water is denser and does not mix with toluene: it collects at
the bottom; only the upper layer reaches the return arm.
\textbf{13.} When water stops collecting, at the expected volume.
\textbf{14.} $0.1000 \times 194.9 = \qty{19.5}{g}$.
\textbf{15.} \ce{C6H11BrO2 + Mg -> C6H11BrMgO2} (dry ether).
\textbf{16.} The reagent adds to \ce{(CH3)2CO}: the \omterm{def:b1:alcohol-activation:alkoxide}{alkoxide}
\ce{(CH3)2C(O^-)CH2CH2CH2-} on the protected chain, with \ce{MgBr+}.
\textbf{17.} Acid would remove the \omterm{def:b1:acetals-protection:protecting-group}{protecting group} too early and could
dehydrate the tertiary alcohol.
\textbf{18.} $0.0700 \times 174.0 = \qty{12.2}{g}$.
\textbf{19.} The \omterm{def:b1:acetals-protection:hemiacetal}{acetal} + \ce{H2O} $\to$ the aldehyde + ethane-1,2-diol
(acid \omterm{def:b1:elementary-steps:catalyst}{catalyst}).
\textbf{20.} Its loss of water needs \omterm{def:b1:predominant-reaction:strong-acid}{strong acid} and heat; mild, cold
aqueous acid hydrolyses the \omterm{def:b1:acetals-protection:hemiacetal}{acetal} much faster.
\textbf{21.} The OH of carbon 5 adds to the aldehyde carbon 1: a six-membered
ring (five carbons and one oxygen) with two methyl groups on the carbon
next to the ring oxygen and an OH on the other carbon next to it.
\textbf{22.} $0.0700 \times 0.90 = \qty{0.0630}{mol}$, \qty{8.19}{g}.
\textbf{23.} $0.0630/0.1000 = \qty{63}{\percent}$ (\omterm{def:b1:acetals-protection:protecting-group}{protection} taken as
complete).
\textbf{24.} $0.1000 \times 18.0 = \qty{1.80}{g}$.
\textbf{25.} $1.80/0.997 =$ \textbf{\qty{1.8}{mL} of water} in the trap.
\end{solution}
